% ============================================
% MEM-06 (EN)
% English translation — SPEC-935-R355
% ============================================
\section*{\textbf{MEM-06} --- The Government Pen}
\markboth{MEM-06 --- The Government Pen}{}

\noindent \textit{The Molambudos --- Concentration Camps of Ceará}

\noindent \textbf{1915 --- Concentration Camps of Ceará}

\vspace{0.5em}\noindent\rule{\textwidth}{0.5pt}\vspace{0.5em}

\noindent The government called them \textquotedblleft{}pens\textquotedblright{}.

\noindent The right word. We were cattle.

\noindent The Alagadiço Camp, on the outskirts of Fortaleza, was an area fenced with barbed wire, about two hundred meters wide and perhaps three hundred meters long. No roof. No walls. Nothing but packed earth, sun, and the smell of people piled up.

\noindent On the first day, I counted the people. I reached three hundred and something before losing count. Then more people arrived. Then more. By the end of the first week, we were squeezed together like animals in transport.

\noindent Every day soldiers arrived with more retirantes. They came from everywhere: Quixeramobim, Boa Viagem, Senador Pompeu, Quixadá, Canindé. The soldiers called us \textquotedblleft{}molambudos\textquotedblright{}. The word came out of their mouths as if it were the name of a plague.

\noindent \textit{\textquotedblleft{}Here come more molambudos.\textquotedblright{}}

\noindent \textit{\textquotedblleft{}Molambudo, get out of there.\textquotedblright{}}

\noindent \textit{\textquotedblleft{}Give food to these molambudos before they all die.\textquotedblright{}}

\vspace{0.5em}\noindent\rule{\textwidth}{0.5pt}\vspace{0.5em}

\noindent Food came once a day.

\noindent A thin broth. The soldiers called it \textquotedblleft{}the government soup\textquotedblright{}. It was water with flour and a few leaves floating in it. I drank mine and still licked the gourd bowl. My mother gave hers to me. I said no. She said she had already eaten. A lie. I could see the cough shaking her empty body.

\noindent The government sent cattle now and then. Not to feed us --- to feed the soldiers. We stood watching them roast the meat. The smell crossed over the fence and entered our stomachs. Some cried. Others beat on the barbed wire. The soldiers laughed.

\noindent \textit{\textquotedblleft{}Let them smell it\textquotedblright{}\textit{, one of them said. }\textquotedblleft{}Hunger is not cured by smell.\textquotedblright{}}

\noindent He was right. Hunger is not cured by smell. Not by crying. Not by prayer. Hunger is cured only by food.

\noindent And there was no food.

\vspace{0.5em}\noindent\rule{\textwidth}{0.5pt}\vspace{0.5em}

\noindent The deaths began in the second week.

\noindent First the old ones. Then the children. Then the adults who were already sick before they arrived.

\noindent The soldiers had a routine. In the morning, they opened the fence, came in, walked among us, and pointed out those who had died during the night. They carried the bodies by the arms and legs --- one on each side --- and took them to the trench.

\noindent The trench lay two hundred yards away. We could see the vultures.

\noindent My mother coughed.

\vspace{0.5em}\noindent\rule{\textwidth}{0.5pt}\vspace{0.5em}

\noindent One night, I dreamed of food.

\noindent It was a good dream. I was at my grandmother's house, who had died before I was born, but in the dream she was alive. She was making corn cake. The cake came out of the pot and rose into the sky. I opened my mouth and it fell into me like rain.

\noindent I woke up hungry.

\noindent A hunger that hurt. Not in the stomach. In the whole body. My bones ached. My mouth ached. My teeth ached. I looked at my hands and saw the bones beneath the skin.

\noindent \textit{\textquotedblleft{}Mother\textquotedblright{}\textit{, I said. }\textquotedblleft{}I want to leave.\textquotedblright{}}

\noindent She did not answer.

\noindent \textit{\textquotedblleft{}Mother?\textquotedblright{}}

\noindent She was lying on her side. Her body was warm. Her breathing was short. Fast. Like a little bird's.

\noindent \textit{\textquotedblleft{}Go\textquotedblright{}\textit{, she said. }\textquotedblleft{}Go, Joaquim. Go.\textquotedblright{}}

\noindent \textit{\textquotedblleft{}Where?\textquotedblright{}}

\noindent She coughed. The cough shook her whole body. When the cough passed, her lips were red.

\noindent \textit{\textquotedblleft{}Far from here.\textquotedblright{}}

\noindent \textit{\textquotedblleft{}I don't know where to go.\textquotedblright{}}

\noindent She closed her eyes.

\noindent \textit{\textquotedblleft{}You know.\textquotedblright{}}

\noindent And I did know.

\noindent I did not want to know. But I knew.

\noindent The trench. The only thing beyond the barbed wire. The only direction that did not require walking. The trench. There was meat there.

\noindent I did not go that night.

\noindent I went the next.\NE{The existence of the Alagadiço Camp, near Fortaleza, during the 1915 drought is documented. Joaquim, his family, and the sensory episodes in this passage are fictional; details of scale, routine, and the trench belong to the narrative reconstruction and are not archival testimony.}

\vspace{0.5em}\noindent\rule{\textwidth}{0.5pt}\vspace{0.5em}

\noindent The Alagadiço Camp was created during the 1915 drought, on the outskirts of Fortaleza. Historical records indicate that thousands of retirantes were confined there, but the sources consulted do not support a precise total of internees or deaths for this narrative.

\noindent I only knew that many did not leave.

\noindent I was one of those who remained.

\noindent Sometimes I think: \textit{why me?}

\noindent Then the creature laughs inside me.

\noindent And I remember that I did not survive \textit{in spite of} the trench.

\noindent I survived \textit{because of} it.

\vspace{0.5em}
\noindent\textit{\textbf{↪ Routes:} \rota{MEM-07} \rotasep \rota{MEM-01}}
